\documentclass{article}

\usepackage{graphicx}
\usepackage{listings}
\usepackage{a4wide}
\usepackage[latin1]{inputenc}
%\usepackage[portuges]{babel}
\usepackage[brazil]{babel}
\usepackage{float}

\title{ 
  \begin{center}
  \includegraphics[scale=0.20]{./imagens/ifal-logo.jpg}
  \end{center}
  \normalsize 
  \textbf{Instituto Federal de Alagoas - IFAL \\ Campus Palmeira do �ndios \\  Programa��o Web \\ Prof. Ivo Calado \\}  
}

\begin{document}

\maketitle \scriptsize
\section{Implementando o Servlet no eclipse}
\begin{enumerate}

\item Criar um projeto Web din�mico no eclipse chamado \emph{ProjetoServlet01};

\item Criar dois pacotes:
\begin{itemize}

\item \emph{modelo}: Essa pasta serve para guardar as
classes/componentes da camada de neg�cio;

\item \emph{web}: Essa pasta serve para guardar os Servlets.

\end{itemize}

\item Implementar a classe Aluno dentro do pacote \emph{modelo},
conforme abaixo:


\begin{lstlisting}
package modelo;

public class Aluno {
    private String nome;

    public Aluno(int codAluno) {
        if (codAluno == 1) {
            this.nome = "Manoel";
        }
        if (codAluno == 2) {
            this.nome = "Pedro";
        }
    }

    public String getNome() {
        return nome;
    }
}
\end{lstlisting}


\item Implementar o primeiro servlet, chamado
\emph{PrimeiroServlet}, conforme abaixo:



\begin{lstlisting}
package web;

import java.io.IOException; 

import javax.servlet.ServletException; 
javax.servlet.http.HttpServlet; 
javax.servlet.http.HttpServletRequest;
javax.servlet.http.HttpServletResponse;

import modelo.Aluno;

public class PrimeiroServlet extends HttpServlet {
    public void doPost(HttpServletRequest request, HttpServletResponse response)
            throws IOException, ServletException {
        response.setContentType("text/html");
        ServletOutputStream out = response.getOutputStream();
        out.println("teste");
    }

    public void doGet(HttpServletRequest request, HttpServletResponse response)
            throws IOException, ServletException {
        response.setContentType("text/html");
        ServletOutputStream out = response.getOutputStream();
        Aluno aluno = new Aluno(1);
        out.println(aluno.getNome());
    }

}
\end{lstlisting}


\item Reiniciar o servi�o do \emph{tomcat};

\item Executar a URl http://localhost:8080/ProjetoServlet01/PrimeiroServlet

\end{enumerate}
\end{document}
